\NativeSlot{44}{60}{%
\ExamGroup{二、选择（将正确答案的序号填在括号里）。{\normalfont\kaishu\fontsize{10.8bp}{14bp}\selectfont （每题2分，共12分）}}

% p106_s01

\ExamQuestion{1.}{5}%
王大爷要给小菜园围上一圈篱笆，\mbox{（\hspace{1.8em}）}的围法更牢固些。\par
\BookFigure{figures/p106_s01-01.tex}
}
\NativeSlot{146}{42}{%
% p106_s02

\ExamQuestion{2.}{5}%
由我国自主研制的大型多用途民用直升机“吉祥鸟”首飞成功，其最大外吊挂起飞质量为\(13.8\)吨，其中的“8”表示\mbox{（\hspace{1.8em}）}。\par
\ExamOptions{A. 8kg}{B. 80kg}{C. 800kg}{D. 8000kg}
}
\NativeSlot{220}{74}{%
% p106_s03

\ExamQuestion{3.}{5}%
下面算式中，计算正确的是\mbox{（\hspace{1.8em}）}。\par
\ExamOptions{A. \(395-78+22=395-(78+22)\)}{B. \(25\times44\times7=25\times4+40\times7\)}{C. \(18\times16\div18\times6=(18\times16)\div(18\times6)\)}{D. \(99\times48=100\times48-48\)}
}
\NativeSlot{344}{77}{%
% p106_s04

\ExamQuestion{4.}{5}%
把一根细铁丝剪两刀分成三段，围成一个三角形（接头处忽略不计），这两刀可以剪在下图中\mbox{（\hspace{1.8em}）}处。\par
\ExamOptions{A. ①②}{B. ②③}{C. ②④}{D. ③④}
\BookFigure{figures/p106_s04-01.tex}
}
\NativeSlot{474}{99}{%
% p106_s05

\ExamQuestion{5.}{5}%
下列图示中，不能表示\(0.04\)的是\mbox{（\hspace{1.8em}）}。\par
\par\noindent%
\ExamTile{.485}{%
A. \raisebox{-21bp}{{\setlength{\NativeFigureHeight}{64bp}\NativeImage{figures/p106_s05-01.tex}}}
}\hfill%
\ExamTile{.485}{%
B. \raisebox{-21bp}{{\setlength{\NativeFigureHeight}{60bp}\NativeImage{figures/p106_s05-02.tex}}}
}\par
\vspace{5bp}
\par\noindent%
\ExamTile{.485}{%
C. \raisebox{-21bp}{{\setlength{\NativeFigureHeight}{59bp}\NativeImage{figures/p106_s05-03.tex}}}
}\hfill%
\ExamTile{.485}{%
D. \raisebox{-21bp}{{\setlength{\NativeFigureHeight}{28bp}\NativeImage{figures/p106_s05-04.tex}}}
}\par
}
\NativeSlot{639}{93}{%
% p106_s06

\ExamQuestion{6.}{5}%
计算图①和图②中小正方形的总个数时，下面图示的计算过程可以解释\mbox{（\hspace{1.8em}）}。\par
\ExamOptions{A. 加法交换律}{B. 乘法交换律}{C. 乘法结合律}{D. 乘法分配律}
\BookFigure{figures/p106_s06-01.tex}
}
